% Documentos da coleta. Texto do formulário fornecido pelo autor (versão final do
% Microsoft Forms), conferido com a exportação de 12 envios (9 a 13/11/2026).
\apendice[ape:formulario]{Formulário aplicado}

Este apêndice reproduz o formulário eletrônico ``Avaliação do Luminthor ---
perfil e usabilidade'', aplicado no Microsoft Forms. Foram omitidos apenas os elementos da interface de edição da
plataforma. As anotações ``Requer resposta'' e ``Opção única'' exibidas pela
plataforma aparecem como ``obrigatória'' e ``opção única''; a numeração de 1 a 20
e os códigos dos itens são os do formulário. A apresentação do aceite (C01) está no
Apêndice~\ref{ape:consentimento}.

A instrução da Seção 2 menciona um código fornecido pelo pesquisador, mas a versão
reproduzida não contém campo para esse código. A numeração dos itens de perfil
também não inclui P01 e P04.

\section*{Apresentação}

\textbf{Avaliação do Luminthor --- perfil e usabilidade}

Versão curta --- perfil do participante e usabilidade do aplicativo.

Para instalar o app acesse: \url{https://luminthor.vercel.app/}

Responda após experimentar o Luminthor. Se não conseguiu instalar ou abrir o
aplicativo, você poderá registrar o problema.

\section*{Seção 1 --- Participação voluntária}

Texto e item C01 reproduzidos no Apêndice~\ref{ape:consentimento}.

\section*{Seção 2 --- Perfil e aparelho}

Use apenas o código fornecido pelo pesquisador. Não escreva seu nome. Se não souber
uma informação técnica, responda ``Não sei''.

\begin{enumerate}
\setcounter{enumi}{1}
\item P02 --- Qual opção descreve melhor sua situação atual? (obrigatória; opção única)
  \begin{itemize}
    \item Estou cursando uma disciplina introdutória de lógica ou programação.
    \item Cursei recentemente uma disciplina introdutória de lógica ou programação.
    \item Estudo programação por conta própria e estou começando.
    \item Já tenho experiência além do nível introdutório.
    \item Ainda não estudei lógica ou programação.
  \end{itemize}
\item P03 --- Há quanto tempo você tem contato com programação? (obrigatória; opção única)
  \begin{itemize}
    \item Ainda não tive contato.
    \item Menos de 3 meses.
    \item De 3 a 6 meses.
    \item De 7 a 12 meses.
    \item Mais de 12 meses.
  \end{itemize}
\item P05 --- Qual é a marca e o modelo do celular utilizado? (obrigatória; texto de
  linha única) Informe apenas marca/modelo, sem número de série ou identificadores.
  ``Não sei'' é uma resposta válida.
\item P06 --- Qual é a versão do Android desse celular? (obrigatória; texto de linha
  única) Você pode consultar as configurações do aparelho. ``Não sei'' é uma resposta
  válida.
\item P07 --- Como você vai consultar o formulário durante os testes? (obrigatória;
  opção única)
  \begin{itemize}
    \item No mesmo celular em que usarei o aplicativo.
    \item Em outro dispositivo, usando o aplicativo no celular.
    \item Em uma versão impressa do formulário.
  \end{itemize}
\end{enumerate}

\section*{Seção 3 --- Instalação e uso do aplicativo}

Instale e abra o APK enviado pelo pesquisador. Na tela inicial, use ``Novo arquivo'';
no editor, ``Execute'' processa o código e o painel inferior mostra saída e
diagnósticos. No depurador, as setas avançam e voltam pelos passos.

Se o aplicativo abrir, faça as três tarefas curtas a seguir e depois responda sobre a
usabilidade. Se houver dificuldade, registre o que conseguiu; não é necessário acertar.

\begin{enumerate}
\setcounter{enumi}{6}
\item I01 --- Você conseguiu instalar e abrir o APK enviado? (obrigatória; opção única)
  Se não conseguiu abrir o aplicativo, você seguirá diretamente para os comentários
  finais.
  \begin{itemize}
    \item Sim, instalei e o aplicativo abriu.
    \item Não consegui instalar.
    \item Instalei, mas o aplicativo não abriu.
    \item Não recebi ou não consegui acessar o APK.
    \item Não tenho um aparelho Android disponível.
    \item Decidi não instalar.
  \end{itemize}
\end{enumerate}

\section*{Seção 4 --- Três tarefas práticas}

Faça as tarefas na ordem, usando um programa por vez. Copie o código no editor,
substituindo o anterior. Tente primeiro sem ajuda externa e use as mensagens do
próprio aplicativo.

Cada tarefa tem um único campo de resposta. Se não conseguir concluir, registre a
dificuldade; se não realizou a tarefa, deixe em branco. Se o aplicativo impedir a
continuação, avance para a avaliação de usabilidade e use ``Não consegui avaliar''
nos recursos não utilizados.

\begin{enumerate}
\setcounter{enumi}{7}
\item T1 --- Executar e acompanhar variáveis (texto multilinha)

Copie e execute o programa. Abra o depurador e avance até o último passo.
\begin{lstlisting}[language={},numbers=none]
var x = 2;

var y = 3;

x = x + y;

y = x * 2;

print(x);
print(y);
\end{lstlisting}
Registre a saída do aplicativo e os valores finais de x e y no depurador. Se teve
dificuldade, descreva-a.

\item T2 --- Corrigir um erro de sintaxe (texto multilinha)

Copie e execute o programa abaixo. Leia o diagnóstico antes de tentar corrigir.
Preserve o cálculo e a impressão do resultado.
\begin{lstlisting}[language={},numbers=none]
var total = 10 + 5;
print(total;
\end{lstlisting}
Em sua resposta, explique o que a mensagem informou, cole o código da sua tentativa
de correção e registre a saída da última execução. Se não conseguiu corrigir,
descreva a dificuldade.

\item T3 --- Investigar uma falha de execução (texto multilinha)

Execute o programa e leia o diagnóstico. Antes de alterar o código, abra o depurador
e observe os valores de total e quantidade no último estado disponível antes da falha.
\begin{lstlisting}[language={},numbers=none]
var total = 12;

var quantidade = 0;

var media = total / quantidade;

print(media);
\end{lstlisting}
Tente corrigir apenas um valor inicial, mantendo o cálculo e a impressão. Registre os
valores observados antes da falha, explique o problema, cole o código da sua tentativa
de correção e informe a saída final. Se não conseguiu, descreva a dificuldade.
\end{enumerate}

\section*{Seção 5 --- Usabilidade do Luminthor}

Responda considerando sua experiência nesta sessão.

Escala: 1 --- Discordo totalmente; 2 --- Discordo; 3 --- Nem concordo nem discordo;
4 --- Concordo; 5 --- Concordo totalmente.

Selecione ``Não consegui avaliar'' quando não tiver usado ou observado o recurso.
Essa opção não é uma nota intermediária.

Os itens 11 a 18 são obrigatórios, de opção única, com as alternativas
``1 --- Discordo totalmente.'', ``2 --- Discordo.'', ``3 --- Nem concordo nem
discordo.'', ``4 --- Concordo.'', ``5 --- Concordo totalmente.'' e ``Não consegui
avaliar.''.

\begin{enumerate}
\setcounter{enumi}{10}
\item U01 --- Foi fácil aprender a usar os controles básicos do Luminthor.
\item U02 --- As mensagens indicaram com clareza onde o erro ocorreu.
\item U03 --- As mensagens explicaram com clareza a causa provável do erro.
\item U04 --- Os termos utilizados nas mensagens foram fáceis de compreender.
\item U05 --- As sugestões apresentadas deixaram claro o próximo passo para corrigir o programa.
\item U06 --- Consegui corrigir os erros sem precisar de explicação adicional.
\item U07 --- O depurador visual ajudou a compreender a mudança dos valores das variáveis.
\item U08 --- A interface foi adequada para escrever e corrigir programas em um smartphone.
\end{enumerate}

\section*{Seção 6 --- Comentários finais}

Os comentários são opcionais. Não inclua nomes ou dados pessoais. Envie apenas uma
resposta por código de participante. Obrigado por registrar sua experiência.

\begin{enumerate}
\setcounter{enumi}{18}
\item A01 --- O que foi mais fácil de usar e o que causou dúvida no Luminthor?
  (texto multilinha)
\item A02 --- Que melhoria você sugere? Se houve problema técnico, descreva o que
  aconteceu. (texto multilinha) Se não conseguiu instalar ou abrir o aplicativo,
  registre aqui a dificuldade. Se recebeu ajuda externa, você também pode
  descrevê-la sem citar nomes.
\end{enumerate}

\apendice[ape:consentimento]{Registro de consentimento}

O texto abaixo corresponde à Seção 1 do formulário aplicado.

\section*{Seção 1 --- Participação voluntária}

Você está sendo convidado(a) a avaliar a usabilidade do Luminthor, aplicativo Android
para programar na linguagem Lumi.

Pesquisador: Vinicius Rodrigues da Costa --- Ciência da Computação, Universidade
Tuiuti do Paraná. Orientador: Prof. Diogenes Cogo Furlan.

A participação é voluntária e não interfere em notas. Você pode desistir e fechar o
formulário antes de enviar. Não é necessário acertar exercícios.

Não informe nome, e-mail, matrícula ou outros dados pessoais. As respostas serão
analisadas para fins acadêmicos e divulgadas em conjunto ou por trechos sem
identificação nominal. O código e os detalhes do relato podem permitir reconhecimento
e serão tratados de forma reservada.

\begin{enumerate}
\item C01 --- Aceita participar voluntariamente desta avaliação? (obrigatória; opção
  única)
  \begin{itemize}
    \item Li as informações e aceito participar.
    \item Não aceito participar.
  \end{itemize}
\end{enumerate}
